\begin{derivation}[从\eqnrefs{eq:sl2c-hermitian-map-dp8}到\eqnrefs{eq:sl2c-double-cover-dp8}]
先计算

\begin{align*}
 \det X
 &=(x^0+x^3)(x^0-x^3)
 -(x^1-\ii x^2)(x^1+\ii x^2)\\
 &=(x^0)^2-(x^1)^2-(x^2)^2-(x^3)^2\\
 &=x_\mu x^\mu.
\end{align*}
对 $A\in SL(2,\mathbb C)$，$X'=AXA^\dagger$ 仍为厄米，且
\begin{align*}
 \det X'
 &=\det A\,\det X\,\det A^\dagger\\
 &=|\det A|^2\det X\\
 &=\det X.
\end{align*}
因此存在唯一实线性变换 $\Phi(A)$ 使 $X(x')=AX(x)A^\dagger$，并且 $\Phi(A)$ 保持 Minkowski 度规。矩阵乘法还给出
\begin{equation*}
 \Phi(A_2A_1)X
 =A_2A_1X A_1^\dagger A_2^\dagger
 =\Phi(A_2)\Phi(A_1)X,
\end{equation*}
所以 $\Phi$ 是群同态。由于 $SL(2,\mathbb C)$ 连通且 $\Phi(\id_2)=\id_4$，其像落在 Lorentz 群的 proper orthochronous 连通分支 $SO^+(1,3)$。

严格求核。若 $A\in\ker\Phi$，则对任意厄米 $X$ 都有
\begin{equation*}
 AXA^\dagger=X.
\end{equation*}
先取 $X=\id_2$，得到 $AA^\dagger=\id_2$，所以 $A$ 是幺正矩阵，$A^\dagger=A^{-1}$。于是上式改写成
\begin{equation*}
 AXA^{-1}=X
 \qquad\text{对任意厄米 }X.
\end{equation*}
厄米矩阵张成整个 $M_2(\mathbb C)$ 复矩阵空间，因此 $A$ 与所有 $2\times2$ 复矩阵对易，只能是标量矩阵 $A=\lambda\id_2$。再用 $\det A=\lambda^2=1$，得到
\begin{equation*}
 \ker\Phi=\{+\id_2,-\id_2\}.
\end{equation*}

最后证明满射，而不是只验证若干特殊变换。把 $\mathfrak{sl}(2,\mathbb C)$ 看成六维实 Lie 代数，其一组实基可取
\begin{equation*}
 -\frac{\ii}{2}\sigma_i,
 \qquad
 +\frac12\sigma_i,
 \qquad i=1,2,3.
\end{equation*}
对转动型无穷小元
\begin{equation*}
 A=\id_2-\frac{\ii}{2}\boldsymbol\theta\cdot\boldsymbol\sigma_{\rm P}
 +O(\theta^2),
\end{equation*}
有
\begin{align*}
 \delta X
 &=-\frac{\ii}{2}
 [\boldsymbol\theta\cdot\boldsymbol\sigma_{\rm P},
   \mathbf x\cdot\boldsymbol\sigma_{\rm P}]\\
 &=(\boldsymbol\theta\times\mathbf x)\cdot\boldsymbol\sigma_{\rm P},
\end{align*}
即 $\delta x^0=0$、$\delta\mathbf x=\boldsymbol\theta\times\mathbf x$，给出三个独立空间转动方向。对推动型无穷小元
\begin{equation*}
 A=\id_2+\frac12\boldsymbol\xi\cdot\boldsymbol\sigma_{\rm P}
 +O(\xi^2),
\end{equation*}
则
\begin{align*}
 \delta X
 &=\frac12\{\boldsymbol\xi\cdot\boldsymbol\sigma_{\rm P},X\}\\
 &=(\boldsymbol\xi\cdot\mathbf x)\id_2
 +x^0\boldsymbol\xi\cdot\boldsymbol\sigma_{\rm P},
\end{align*}
因此 $\delta x^0=\boldsymbol\xi\cdot\mathbf x$、$\delta\mathbf x=x^0\boldsymbol\xi$，给出三个独立 Lorentz 推动方向。故微分
\begin{equation*}
 \dd\Phi_e:\mathfrak{sl}(2,\mathbb C)_{\mathbb R}
 \longrightarrow\mathfrak{so}(1,3)
\end{equation*}
在六个实维度上满秩，因而是 Lie 代数同构。于是 $\Phi(SL(2,\mathbb C))$ 是 $SO^+(1,3)$ 的开子群；而 $SO^+(1,3)$ 连通，开子群的不同陪集若多于一个会把群分成互不相交的非空开集，所以只能有一个陪集。于是 $\Phi$ 满射。

综上，$\Phi$ 是满射群同态，核为 $\{\pm\id_2\}$，由第一同构定理得到
\begin{equation*}
 SL(2,\mathbb C)/\{\pm\id_2\}
 \cong SO^+(1,3),
\end{equation*}
即正文\eqnrefs{eq:sl2c-double-cover-dp8}。这里的“双覆盖”因此同时包含“恰有二元核”和“确实覆盖整个 $SO^+(1,3)$”两层结论。
\end{derivation}
